% !TeX program = lualatex
% ============================================================
% chapter7.tex — ch07 唯一內容源（2026-08-09 起；LaTeX 統一 U1）
% 沿革：升格自 dist/ch07/chapter7.tex（原 make_dist.py 自 fragment 轉換的成品；
%   fragment 樹已凍結，改課文一律改本檔。拍板見 ../../KICKOFF-latex-unification.md）
%   樣式源＝../../template/calcbook.sty（M-B1 拍板模板）
% 編譯（本資料夾為 CWD；aux 進 build/、成品 PDF 移入 ../../dist/ch07/）：
%   latexmk -lualatex -auxdir=../../build/aux-ch07 chapter7.tex
% ============================================================
\documentclass[a4paper,12pt,oneside]{memoir}
\makeatletter\def\input@path{{../../template/}}\makeatother
\def\cbfontsdir{../../template/fonts/inter/}
\usepackage{calcbook}
% 圖：emitter 射 `<ch>/<stem>`（見 convert.py），故 graphicspath 指向 chapters/ 上層。
% 模板內的 {../figs/} 是 chapters/<ch>/ 為 CWD 時的路徑，dist/<ch>/ 需這一行覆寫。
\graphicspath{{../../chapters/}}
\begin{document}
\cbchapter{7}
\chapteropener{Chapter 7}{Applications of Integration}
\begin{lead}
Chapter 6 built the machine. The definite integral was assembled from Riemann sums, and the Fundamental Theorem of Calculus turned its evaluation into a search for antiderivatives. This chapter puts the machine to work on the area between two curves, the volume of a solid, the work done by a varying force, the average of a continuously varying quantity, the length of a curve, and the area of a curved surface. Six problems, and one method: \emph{slice} the quantity into thin pieces, \emph{approximate} each piece by a product, \emph{sum} the pieces, and \emph{refine}. The product is a value that is nearly constant across the slice, times the slice’s width. Every time, the sums are Riemann sums, their limit is a definite integral, and the Fundamental Theorem evaluates it. What is genuinely new in this chapter is therefore not the integration but the setup: deciding what to slice, reading off what each thin slice contributes, and recognizing the resulting limit as an integral we can compute. The skill this chapter teaches is the modelling step of representing a quantity by an integral, a skill the rest of calculus keeps drawing on, from differential equations to the geometry of space.
\end{lead}
\parahead{By the end of this chapter, you will be able to:}
\begin{objectives}
  \item compute the area of a region bounded by curves, slicing along either axis;
  \item compute volumes by cross-sections, by disks and washers, and by cylindrical shells;
  \item set up and evaluate work integrals for springs, cables, and pumping problems;
  \item compute the average value of a function and locate a point where it is attained;
  \item compute the length of a smooth curve and the area of a surface of revolution, and work with the arc length function.
\end{objectives}
\sechead{7.1}{Areas Between Curves}
Chapter 6 measured the region between a curve and the \(x\)-axis. Most regions are not like that: they are trapped between \emph{two} curves, a curved roof and a curved floor. The chapter’s slicing schema handles them with one adjustment. Cut the region into thin vertical strips. Each strip runs from the lower curve up to the upper one, so its height is not the height of either graph above the axis but the \emph{gap} between the graphs. Approximate the strip by a rectangle, gap times width; sum; refine. The axis has dropped out of the computation entirely, and that is precisely what makes the method robust.

\subsechead{The area between two curves}
Let \(f\) and \(g\) be continuous on \([a, b]\) with \(f(x) \ge g(x)\) for all \(x\) in \([a, b]\), and let \(S\) be the region bounded above by \(y = f(x)\), below by \(y = g(x)\), and on the sides by the vertical lines \(x = a\) and \(x = b\). Divide \([a, b]\) into \(n\) strips of equal width \(\Delta x = \tfrac{b-a}{n}\), and in the \(i\)-th strip pick a sample point \(x_{i}^{*}\). The strip is approximated by a rectangle of width \(\Delta x\) and height \(f(x_{i}^{*}) - g(x_{i}^{*})\), so the region’s area is approximated by the sum

\[ \sum_{i=1}^{n} \bigl[f(x_{i}^{*}) - g(x_{i}^{*})\bigr]\,\Delta x. \]

This is a Riemann sum, not for \(f\) or for \(g\), but for the single function \(f - g\). Since \(f - g\) is continuous on \([a, b]\), Theorem 6.1 guarantees that as \(n \to \infty\) these sums converge to \(\int_{a}^{b} \bigl[f(x) - g(x)\bigr]\,dx\), no matter how the sample points are chosen. As in §6.1, we take the common limit as the \emph{definition} of the area.

\begin{figureblock}
  \includegraphics[width=74.79mm]{ch07/figs/gap-strip}
\figcaption{fig:7.1}{The region between an upper curve and a lower one, cut into vertical strips. The strip’s height is the gap \(f(x_i^{*}) - g(x_i^{*})\) between the graphs, not a height above the axis, which passes through the region without playing any role.}
\end{figureblock}
\begin{envdefinition}{Definition}{def:7.1}{Area between curves}
Let \(f\) and \(g\) be continuous on \([a, b]\) with \(f(x) \ge g(x)\) on \([a, b]\). The \emph{area} of the region bounded by \(y = f(x)\) above, \(y = g(x)\) below, and the lines \(x = a\) and \(x = b\) is

\[ A = \int_{a}^{b} \bigl[f(x) - g(x)\bigr]\,dx. \]
\end{envdefinition}
Two checks tie the definition to what we already have. First, when \(g\) is the zero function and \(f \ge 0\), the region is the one under the graph of \(f\), and the formula collapses to \(\int_{a}^{b} f(x)\,dx\), the area formula of Definition 6.1. Second, the definition asks only that \(f\) stay \emph{above} \(g\); neither curve is required to stay above the axis.

That second point deserves a moment, because in §6.2 we saw that below the axis an integral counts area \emph{negatively}. Why does no such caution attach here? Raise both curves by the same constant \(C\). That is, replace \(f\) by \(f + C\) and \(g\) by \(g + C\) with \(C\) large enough to lift the whole region above the axis. The region merely slides upward, so its area is unchanged; and the integrand is unchanged too, since \((f + C) - (g + C) = f - g\). The subtraction has already discarded the only thing the axis was ever needed for. Subtract the graphs first, and signs take care of themselves.

\begin{workedexample}
\begin{envexample}{Example}{ex:7.1}{}
Find the area of the region bounded above by \(y = e^{x}\), below by \(y = x\), and on the sides by \(x = 0\) and \(x = 1\).
\end{envexample}
\begin{envsolution}{Solution}{}{}
First confirm which curve is on top throughout the interval: \(e^{x}\) is increasing, so \(e^{x} \ge e^{0} = 1\) on \([0, 1]\), while \(x \le 1\) there; hence \(e^{x} \ge 1 \ge x\) across the whole interval. Definition \ref{def:7.1} applies with \(f(x) = e^{x}\) and \(g(x) = x\), and the Fundamental Theorem (Theorem 6.4) evaluates the result:

\[ \begin{aligned}
          A &= \int_{0}^{1} \bigl(e^{x} - x\bigr)\,dx = \left[e^{x} - \frac{x^{2}}{2}\right]_{0}^{1} \\[2pt]
            &= \left(e - \frac{1}{2}\right) - (1 - 0) = e - \frac{3}{2} \approx 1.22.
        \end{aligned} \]\qedmark
\end{envsolution}
\end{workedexample}
In Example \ref{ex:7.1} the interval \([a, b]\) was handed to us. More often the region is \emph{enclosed} by the curves themselves, and the first task is to find where they meet: the intersection points supply the limits of integration.

\begin{workedexample}
\begin{envexample}{Example}{ex:7.2}{}
Find the area of the region enclosed by the parabolas \(y = x^{2}\) and \(y = 2x - x^{2}\).
\end{envexample}
\begin{envsolution}{Solution}{}{}
The curves intersect where \(x^{2} = 2x - x^{2}\), that is, \(2x^{2} - 2x = 0\), so \(2x(x - 1) = 0\) and \(x = 0\) or \(x = 1\). The enclosed region lives over the interval \([0, 1]\). To see which parabola is on top there, test a point between the intersections: at \(x = \tfrac{1}{2}\), the first curve gives \(\tfrac{1}{4}\) and the second gives \(1 - \tfrac{1}{4} = \tfrac{3}{4}\). So \(y = 2x - x^{2}\) runs above \(y = x^{2}\), and

\[ \begin{aligned}
          A &= \int_{0}^{1} \bigl[(2x - x^{2}) - x^{2}\bigr]\,dx = \int_{0}^{1} \bigl(2x - 2x^{2}\bigr)\,dx \\[2pt]
            &= \left[x^{2} - \frac{2x^{3}}{3}\right]_{0}^{1} = 1 - \frac{2}{3} = \frac{1}{3}.
        \end{aligned} \]\qedmark
\end{envsolution}
\end{workedexample}
\begin{figureblock}
  \includegraphics[width=70.64mm]{ch07/figs/enclosed-lens}
\figcaption{fig:7.2}{The region enclosed by \(y = x^{2}\) and \(y = 2x - x^{2}\). No vertical sides are given or needed: the two intersection points close the region by themselves, and they become the limits of integration.}
\end{figureblock}
\subsechead{When the curves cross}
Definition \ref{def:7.1} assumes one curve stays on top. When the curves cross inside the interval, the roles of roof and floor swap at each crossing, and the honest quantity to track is the size of the gap, \(\lvert f(x) - g(x)\rvert\). The total area between the graphs is

\[ A = \int_{a}^{b} \bigl\lvert f(x) - g(x)\bigr\rvert\,dx, \]

and when the crossings are finitely many — as they are in every example we will meet — we compute it by splitting the interval at the crossings and applying Definition \ref{def:7.1} on each piece with the curves in their correct local order. This is the same move §6.4 made for total distance: displacement integrates \(v\), but distance travelled integrates \(\lvert v\rvert\), splitting where the velocity changes sign.

\begin{envcaution}{Caution}{}{}
If the curves cross inside \([a, b]\), the single integral \(\int_{a}^{b} (f - g)\,dx\) adds \emph{signed} gaps: pieces where \(g\) is on top count negatively and cancel against the rest, exactly as signed area cancelled in §6.2. An answer that comes out negative, or suspiciously small, usually means a crossing was missed. Find every crossing first; then split, or integrate \(\lvert f - g\rvert\).
\end{envcaution}
\begin{workedexample}
\begin{envexample}{Example}{ex:7.3}{}
Find the total area of the region between \(y = \sin x\) and \(y = \cos x\) for \(0 \le x \le \pi/2\).
\end{envexample}
\begin{envsolution}{Solution}{}{}
The curves cross where \(\sin x = \cos x\). On \([0, \pi/2)\) the cosine is positive, so we may divide by it: \(\tan x = 1\), giving \(x = \pi/4\); and \(x = \pi/2\) is not a crossing since \(\sin\tfrac{\pi}{2} = 1 \ne 0 = \cos\tfrac{\pi}{2}\). At \(x = 0\) the cosine is on top (\(1 > 0\)), so \(\cos x \ge \sin x\) on \([0, \pi/4]\) and the order swaps on \([\pi/4, \pi/2]\). Split there:

\[ \begin{aligned}
          A &= \int_{0}^{\pi/4} (\cos x - \sin x)\,dx + \int_{\pi/4}^{\pi/2} (\sin x - \cos x)\,dx \\[2pt]
            &= \bigl[\sin x + \cos x\bigr]_{0}^{\pi/4} + \bigl[-\cos x - \sin x\bigr]_{\pi/4}^{\pi/2} \\[2pt]
            &= \left(\tfrac{\sqrt{2}}{2} + \tfrac{\sqrt{2}}{2}\right) - (0 + 1) + (-0 - 1) + \left(\tfrac{\sqrt{2}}{2} + \tfrac{\sqrt{2}}{2}\right) \\[2pt]
            &= \bigl(\sqrt{2} - 1\bigr) + \bigl(\sqrt{2} - 1\bigr) = 2\sqrt{2} - 2 \approx 0.83.
        \end{aligned} \]

The two pieces came out equal, and that is no accident: the identity \(\cos x = \sin\bigl(\tfrac{\pi}{2} - x\bigr)\) reflects the second region onto the first across the line \(x = \pi/4\). Symmetries like this are worth a glance before integrating. Here, computing one piece and doubling it would have been enough. \qedmark
\end{envsolution}
\end{workedexample}
\begin{figureblock}
  \includegraphics[width=75.93mm]{ch07/figs/sincos-lobes}
\figcaption{fig:7.3}{The two lobes between \(y = \cos x\) and \(y = \sin x\) on \([0, \pi/2]\). The roles of roof and floor swap at the crossing \(x = \pi/4\), which is why one fixed order of subtraction would cancel; the identity \(\cos x = \sin\bigl(\tfrac{\pi}{2} - x\bigr)\) makes the second lobe the mirror image of the first.}
\end{figureblock}
\subsechead{Slicing horizontally}
Nothing privileges vertical strips. Suppose a region has a left boundary \(x = g(y)\) and a right boundary \(x = f(y)\), with \(f\) and \(g\) continuous on \([c, d]\) and \(f(y) \ge g(y)\) there. Curves like these read more naturally as \(x\) in terms of \(y\), so slice the region horizontally instead. The strip at height \(y\) runs from the left curve to the right one, its length is the gap \(f(y) - g(y)\), and the same limit argument, resting on the same continuity, gives

\[ A = \int_{c}^{d} \bigl[f(y) - g(y)\bigr]\,dy, \]

with \emph{right minus left} in place of top minus bottom. Choosing the slicing direction well can spare real work, as the next example shows.

\begin{workedexample}
\begin{envexample}{Example}{ex:7.4}{}
Find the area of the region enclosed by the parabolas \(x = y^{2}\) and \(x = 2 - y^{2}\).
\end{envexample}
\begin{envsolution}{Solution}{}{}
Both curves give \(x\) directly in terms of \(y\): the first opens rightward from the origin, the second opens leftward from \((2, 0)\). They meet where \(y^{2} = 2 - y^{2}\), so \(y^{2} = 1\) and \(y = \pm 1\), giving the points \((1, 1)\) and \((1, -1)\). Vertical strips would be awkward here: below \(x = 1\) a strip runs between the two branches of \(x = y^{2}\), beyond it between the branches of \(x = 2 - y^{2}\), forcing a split at \(x = 1\) and square roots throughout. A horizontal strip at height \(y\), by contrast, always enters the region on \(x = y^{2}\) and leaves it on \(x = 2 - y^{2}\). With right minus left,

\[ A = \int_{-1}^{1} \bigl[(2 - y^{2}) - y^{2}\bigr]\,dy = \int_{-1}^{1} \bigl(2 - 2y^{2}\bigr)\,dy. \]

The integrand is even, so by Theorem 6.8 the integral is twice its value over \([0, 1]\):

\[ A = 2\int_{0}^{1} \bigl(2 - 2y^{2}\bigr)\,dy = 2\left[2y - \frac{2y^{3}}{3}\right]_{0}^{1} = 2\left(2 - \frac{2}{3}\right) = \frac{8}{3}. \]\qedmark
\end{envsolution}
\end{workedexample}
\begin{figureblock}
  \includegraphics[width=75.93mm]{ch07/figs/horizontal-slice}
\figcaption{fig:7.4}{The region enclosed by \(x = y^{2}\) and \(x = 2 - y^{2}\), sliced horizontally. Every horizontal strip enters on one parabola and leaves on the other, whereas a vertical strip would change character at \(x = 1\).}
\end{figureblock}
The gap between two curves need not be a geometric height at all. Suppose two runners leave the same start together, so that \(s_{A}(0) = s_{B}(0)\). They run with continuous velocities \(v_{A} = s_{A}'\) and \(v_{B} = s_{B}'\), with \(v_{A}(t) \ge v_{B}(t)\) over the first \(T\) seconds. Each thin time-slice contributes \(\bigl[v_{A}(t) - v_{B}(t)\bigr]\,\Delta t\), the little bit of extra ground the faster runner gains during that instant. Since \(s_{A} - s_{B}\) has continuous derivative \(v_{A} - v_{B}\), the Net Change Theorem (Theorem 6.5) gives \(\int_{0}^{T} \bigl[v_{A}(t) - v_{B}(t)\bigr]\,dt = s_{A}(T) - s_{B}(T)\), the faster runner’s lead at time \(T\). On a shared velocity–time graph, that lead \emph{is} the area between the two velocity curves. The same reading works for any two rates, such as two flow rates or two growth rates, \emph{as long as one of them stays at least as large as the other}. If the rate curves cross, the area between the graphs accumulates the absolute gap \(\int \lvert v_{A} - v_{B}\rvert\), while the net difference is the plain integral. That is the displacement-versus-distance distinction of §6.4 again.

\begin{envstrategy}{Strategy}{strat:7.1}{Setting up an area integral}
\begin{steps}
  \item \emph{Find the boundaries.} Sketch the curves and find every intersection in the range of interest by solving \(f = g\); intersections supply the limits of integration when the region is enclosed.
  \item \emph{Choose the slicing direction.} Slice so that each strip meets the same two boundary curves throughout: vertically when the boundaries are graphs \(y = f(x)\), \(y = g(x)\); horizontally when they read naturally as \(x\) in terms of \(y\).
  \item \emph{Integrate the gap.} The integrand is top minus bottom (vertical strips) or right minus left (horizontal strips), never a height above an axis.
  \item \emph{Split at crossings.} If the boundary curves swap roles inside the interval, split there; the total area is \(\int \lvert f - g\rvert\), computed piece by piece.
\end{steps}
\end{envstrategy}
One move carried this whole section. Once we subtract the boundary graphs, the machinery of Chapter 6 applies to the difference exactly as it stood: Riemann sums, the existence guarantee of Theorem 6.1, evaluation by the Fundamental Theorem. The area between curves is the integral of a gap, and the runners’ race shows the same reading paying off outside geometry. The next section keeps the slicing and changes the dimension: the slices become thin solid slabs, the gap becomes a cross-sectional area, and the same sum-and-refine limit measures volume.


\sechead{7.2}{Volumes}
From flat regions we turn to solids, and the slices are now cut from a solid body, the way a loaf of bread is cut into slices. What does geometry give us to work with? For areas it gave rectangles. For volumes it gives the \emph{cylinder}: a solid with a flat base of area \(A\) swept straight up to height \(h\) has volume \(V = Ah\). (The base need not be a circle. A box is a cylinder over a rectangle.) A curved solid is not a cylinder anywhere; but each \emph{thin} slice of it very nearly is, and that is all the slicing method needs.

\subsechead{Volume by cross-sections}
Let \(S\) be a solid lying between the vertical planes \(x = a\) and \(x = b\). For each \(x\) in \([a, b]\), slice \(S\) by the plane through \(x\) perpendicular to the \(x\)-axis, and write \(A(x)\) for the area of the exposed face, called the \emph{cross-section} at \(x\). Now cut \([a, b]\) into \(n\) subintervals of width \(\Delta x\) and, in the \(i\)-th one, pick a sample point \(x_{i}^{*}\). The slab of \(S\) over that subinterval is nearly a cylinder with base area \(A(x_{i}^{*})\) and thickness \(\Delta x\), so the whole solid is approximated by

\[ V \approx \sum_{i=1}^{n} A(x_{i}^{*})\,\Delta x, \]

a Riemann sum for the function \(A\). If \(A\) is continuous, Theorem 6.1 guarantees these sums close on a single limit, independent of every choice made in the slicing. It is that limit we take the volume to \emph{be}. Note the word: this is a definition, not a theorem. Geometry gives us the volume of a cylinder. For a curved solid, the modelling step is to define volume as the limit of the volumes of its slab approximations, a step we take once and state openly, just as §6.1 did for area.

\begin{envdefinition}{Definition}{def:7.2}{Volume by cross-sections}
Let \(S\) be a solid lying between the planes \(x = a\) and \(x = b\), and suppose its cross-section at \(x\), perpendicular to the \(x\)-axis, has area \(A(x)\), where \(A\) is continuous on \([a, b]\). The \emph{volume} of \(S\) is

\[ V = \int_{a}^{b} A(x)\,dx. \]
\end{envdefinition}
The definition passes its first test at once. Suppose the cross-section never changes, so that \(A(x) = A\) for all \(x\), as in an actual cylinder lying on its side with height \(h = b - a\). Then \(V = \int_{a}^{b} A\,dx = A\,(b - a) = Ah\), the formula we started from. The examples below run harder tests: the definition will be asked to reproduce two volumes that classical geometry already knows.

\begin{figureblock}
  \includegraphics[width=79.9mm]{ch07/figs/slab-face}
\figcaption{fig:7.5}{A solid sliced like a loaf: the slab at \(x_i^{*}\) is nearly a cylinder, a flat face of area \(A(x_i^{*})\) swept through the small thickness \(\Delta x\), here exaggerated for visibility. The solid is deliberately irregular: nothing in Definition \ref{def:7.2} asks for symmetry, only for a measurable cross-section at each station between \(a\) and \(b\).}
\end{figureblock}
The idea is older than the integral. In the seventeenth century Cavalieri argued that two solids whose cross-sections agree at every height must enclose the same volume. This is \emph{Cavalieri’s principle}. Definition \ref{def:7.2} is that principle made quantitative: the volume is determined by, and computed from, the cross-section areas alone.

\subsechead{Solids of revolution: disks}
The richest supply of solids with computable cross-sections comes from \emph{revolution}. Take the region under \(y = f(x)\) from \(a\) to \(b\) (with \(f\) continuous and \(f \ge 0\)) and rotate it about the \(x\)-axis. Each vertical strip of the region sweeps out a thin circular slab, called a \emph{disk}. The cross-section of the resulting solid at \(x\) is a full circle whose radius is the height of the curve:

\[ A(x) = \pi \bigl[f(x)\bigr]^{2}. \]

Definition \ref{def:7.2} then gives \(V = \int_{a}^{b} \pi \bigl[f(x)\bigr]^{2}\,dx\). No new formula is being introduced here, only a recipe for \(A(x)\).

\begin{envcaution}{Caution}{}{}
For rotation about the \(x\)-axis, the disk’s radius at \(x\) is the \emph{function value} \(f(x)\), the vertical distance from the axis to the curve. The rotation sweeps that segment around the axis. It is not a distance along the \(x\)-axis, and when the axis of rotation moves (as it will shortly), the radius is always measured from \emph{that} axis, not from a coordinate axis out of habit.
\end{envcaution}
\begin{figureblock}
  \includegraphics[width=60.06mm]{ch07/figs/strip-to-disk-1}\hspace{6mm}\includegraphics[width=60.06mm]{ch07/figs/strip-to-disk-2}
\figcaption{fig:7.6}{The same station twice: a vertical strip of height \(f(x)\), and the thin disk it sweeps out when the region rotates about the \(x\)-axis. The disk’s radius is the function value \(f(x)\), its centre sits on the axis, and the solid extends symmetrically below the axis.}
\end{figureblock}
\begin{workedexample}
\begin{envexample}{Example}{ex:7.5}{}
Show that a sphere of radius \(r\) has volume \(\dfrac{4}{3}\pi r^{3}\).
\end{envexample}
\begin{envsolution}{Solution}{}{}
Generate the sphere by rotating the upper half-disk under \(y = \sqrt{r^{2} - x^{2}}\), \(-r \le x \le r\), about the \(x\)-axis. The cross-section at \(x\) is a disk of radius \(\sqrt{r^{2} - x^{2}}\), so

\[ A(x) = \pi \bigl(\sqrt{r^{2} - x^{2}}\bigr)^{2} = \pi \bigl(r^{2} - x^{2}\bigr), \]

which is continuous, and Definition \ref{def:7.2} applies. The integrand is even, so Theorem 6.8 lets us integrate over \([0, r]\) and double:

\[ \begin{aligned}
          V &= \int_{-r}^{r} \pi \bigl(r^{2} - x^{2}\bigr)\,dx = 2\pi \int_{0}^{r} \bigl(r^{2} - x^{2}\bigr)\,dx \\[2pt]
            &= 2\pi \left[r^{2}x - \frac{x^{3}}{3}\right]_{0}^{r} = 2\pi\cdot\frac{2r^{3}}{3} = \frac{4}{3}\pi r^{3}.
        \end{aligned} \]

This is the volume formula the Greeks knew. Definition \ref{def:7.2} has reproduced a fact of classical geometry from nothing but slicing, the same reassurance §6.2 drew from the semicircle in Example 6.6. A definition that agrees with everything already known deserves some confidence. \qedmark
\end{envsolution}
\end{workedexample}
\begin{workedexample}
\begin{envexample}{Example}{ex:7.6}{}
The region under \(y = \sqrt{x}\) from \(0\) to \(4\) is rotated about the \(x\)-axis. Find the volume of the resulting solid.
\end{envexample}
\begin{envsolution}{Solution}{}{}
The cross-section at \(x\) is a disk of radius \(\sqrt{x}\), with area \(A(x) = \pi (\sqrt{x})^{2} = \pi x\), continuous on \([0, 4]\). Hence

\[ V = \int_{0}^{4} \pi x\,dx = \pi \left[\frac{x^{2}}{2}\right]_{0}^{4} = 8\pi. \]

Squaring the radius \emph{before} integrating is what makes this integral easy: the square root disappears into the disk area. \qedmark
\end{envsolution}
\end{workedexample}
\subsechead{Washers}
Rotate the region \emph{between} two curves and the solid acquires a hole. If \(f\) and \(g\) are continuous on \([a, b]\) with \(f(x) \ge g(x) \ge 0\), and the region between the graphs is rotated about the \(x\)-axis, the cross-section at \(x\) is an annulus with outer radius \(R = f(x)\) and inner radius \(r = g(x)\). We call such a ring a \emph{washer}. Its area is a difference of two disk areas,

\[ A(x) = \pi \bigl[f(x)\bigr]^{2} - \pi \bigl[g(x)\bigr]^{2} = \pi \Bigl(\bigl[f(x)\bigr]^{2} - \bigl[g(x)\bigr]^{2}\Bigr), \]

and by the linearity of the integral (Theorem 6.2) the volume splits the same way: the solid with the hole is the full solid minus the hole’s own volume. Where §7.1 subtracted \emph{heights}, the washer subtracts \emph{squared radii}, and the order of operations there is a genuine trap.

\begin{envcaution}{Caution}{}{}
A washer’s area is \(\pi\bigl(R^{2} - r^{2}\bigr)\), \emph{not} \(\pi (R - r)^{2}\): square each radius first, then subtract. Writing \(\pi\bigl(f(x) - g(x)\bigr)^{2}\) squares the gap, as if §7.1’s integrand could be recycled. It is the most common volume error there is, and it is a silent one. The integral still computes, just to the wrong number.
\end{envcaution}
\begin{figureblock}
  \includegraphics[width=62.71mm]{ch07/figs/washer-annulus}
\figcaption{fig:7.7}{A washer face-on: the ring between an outer circle of radius \(R\) and an inner one of radius \(r\), both measured from the axis of rotation at the centre. Its area is the big disk minus the small one, \(\pi\bigl(R^{2} - r^{2}\bigr)\), visibly not the disk of radius \(R - r\).}
\end{figureblock}
\begin{workedexample}
\begin{envexample}{Example}{ex:7.7}{}
The region enclosed by \(y = x\) and \(y = x^{2}\) is rotated about the \(x\)-axis. Find the volume of the resulting solid.
\end{envexample}
\begin{envsolution}{Solution}{}{}
The curves meet where \(x = x^{2}\), so \(x = 0\) or \(x = 1\), and between these values \(x \ge x^{2} \ge 0\) (at \(x = \tfrac12\): \(\tfrac12\) against \(\tfrac14\)). Rotating, the line is the outer boundary and the parabola the inner one: \(R = x\), \(r = x^{2}\), giving washers of area \(\pi\bigl(x^{2} - x^{4}\bigr)\). Hence

\[ V = \int_{0}^{1} \pi \bigl(x^{2} - x^{4}\bigr)\,dx = \pi \left[\frac{x^{3}}{3} - \frac{x^{5}}{5}\right]_{0}^{1} = \pi \left(\frac{1}{3} - \frac{1}{5}\right) = \frac{2\pi}{15}. \]

Had the same region been rotated about the line \(y = 2\) instead, only where the radii are measured from would change. Each radius becomes a distance to \emph{that} line, not to a coordinate axis out of habit. The next example carries that change through. \qedmark
\end{envsolution}
\end{workedexample}
\begin{workedexample}
\begin{envexample}{Example}{ex:7.8}{}
The same region — enclosed by \(y = x\) and \(y = x^{2}\) — is now rotated about the horizontal line \(y = 2\). Find the volume of the resulting solid.
\end{envexample}
\begin{envsolution}{Solution}{}{}
The region still runs over \([0, 1]\) with \(y = x\) on top and \(y = x^{2}\) below, and it lies entirely beneath the axis \(y = 2\). Rotating about that line, each radius is a \emph{vertical distance to \(y = 2\)}. The boundary farther from the axis is now the lower curve \(y = x^{2}\), so it supplies the outer radius \(R = 2 - x^{2}\); the nearer boundary \(y = x\) supplies the inner radius \(r = 2 - x\). The washer at \(x\) has area

\[ \pi\bigl(R^{2} - r^{2}\bigr) = \pi\bigl[(2 - x^{2})^{2} - (2 - x)^{2}\bigr] = \pi\bigl(x^{4} - 5x^{2} + 4x\bigr), \]

so

\[ \begin{aligned}
          V &= \pi\int_{0}^{1} \bigl(x^{4} - 5x^{2} + 4x\bigr)\,dx = \pi\left[\frac{x^{5}}{5} - \frac{5x^{3}}{3} + 2x^{2}\right]_{0}^{1} \\[2pt]
            &= \pi\left(\frac{1}{5} - \frac{5}{3} + 2\right) = \frac{8\pi}{15}.
        \end{aligned} \]

Rotating the same region about the \(x\)-axis gave \(\tfrac{2\pi}{15}\) (Example \ref{ex:7.7}); about \(y = 2\) the solid is different, and larger. The only thing that changed was where the radii were measured from, exactly the discipline the Caution insisted on. \qedmark
\end{envsolution}
\end{workedexample}
\subsechead{Cross-sections without revolution}
Definition \ref{def:7.2} never mentions revolution. It applies to any solid whose cross-sectional area we can describe, and for many classical solids the description is a matter of similar triangles.

\begin{workedexample}
\begin{envexample}{Example}{ex:7.9}{}
Find the volume of a pyramid whose base is a square of side \(b\) and whose height is \(h\).
\end{envexample}
\begin{envsolution}{Solution}{}{}
Place the apex at the origin with the pyramid’s central axis along the \(x\)-axis, so the base sits in the plane \(x = h\). The cross-section at \(x\) is a square; write \(s(x)\) for its side. Slicing the pyramid through its axis exposes a triangle with apex at the origin, and for \(0 < x \le h\) the cross-section’s half-width at distance \(x\) relates to the base’s half-width by similar triangles:

\[ \frac{s(x)/2}{x} = \frac{b/2}{h}, \qquad\text{so}\qquad s(x) = \frac{b}{h}\,x. \]

At the apex itself \(s(0) = 0\), which the same formula reports, so \(s(x) = \tfrac{b}{h}x\) holds across all of \([0, h]\). Hence \(A(x) = s(x)^{2} = \dfrac{b^{2}}{h^{2}}\,x^{2}\), continuous on \([0, h]\), and

\[ V = \int_{0}^{h} \frac{b^{2}}{h^{2}}\,x^{2}\,dx = \frac{b^{2}}{h^{2}}\cdot\frac{h^{3}}{3} = \frac{1}{3}\,b^{2}h. \]

This is the classical one-third rule for pyramids. The integral also shows where the mysterious \(\tfrac13\) comes from: it is \(\int_{0}^{1} u^{2}\,du\), the price of cross-sections that grow like the \emph{square} of the distance from the apex. \qedmark
\end{envsolution}
\end{workedexample}
\begin{figureblock}
  \includegraphics[width=75.93mm]{ch07/figs/pyramid-similar}
\figcaption{fig:7.8}{The pyramid sliced through its central axis: apex at the origin, base at \(x = h\). The exposed triangle carries the whole solution. The half-widths \(s(x)/2\) and \(b/2\) stand over \(x\) and \(h\) in the same ratio, and the inset shows the square cross-section of side \(s(x)\) face-on.}
\end{figureblock}
The section’s whole content is one definition and a habit. The definition: volume is the integral of cross-sectional area, the slicing schema with the strip’s gap promoted to a slab’s face. The habit: before integrating anything, describe the cross-section — disk, washer, or square — and let its geometry write the integrand. The definition justified that confidence by reproducing the cylinder, the sphere, and the pyramid of classical geometry. But perpendicular slices are not always the simplest choice. For some solids of revolution every washer requires solving for \(x\) in terms of \(y\), whereas decomposing the solid into thin cylindrical \emph{shells} gives the needed dimensions directly. That method is the next section’s business.


\sechead{7.3}{Volumes by Cylindrical Shells}
Try to compute a volume with the last section’s tools and you will sometimes find that those methods do not fit the region. Rotate the region under \(y = 2x^{2} - x^{3}\), from \(x = 0\) to \(x = 2\), about the \(y\)-axis. Cross-sections perpendicular to the axis of rotation are washers, so we need the region described from the \(y\)-axis’s point of view: for each height \(y\), which \(x\)-values lie in the region? That means solving \(y = 2x^{2} - x^{3}\) for \(x\), which requires inverting a cubic. That inversion has no simple answer. Yet the region itself could hardly be simpler: read \emph{vertically}, it is just the strips under a graph. The lesson is that slicing is a \emph{choice}, and the right choice is the one whose slices see the region simply.

This section develops the other natural choice for solids of revolution: keep the vertical strip — the one that reads \(f\) directly — and let the rotation carry it \emph{around} the axis. The strip sweeps out a thin-walled tube, a \emph{cylindrical shell}, and the solid is built of nested shells the way a tree trunk is built of rings.

\begin{figureblock}
  \includegraphics[width=52.35mm]{ch07/figs/shell-motivation}
\figcaption{fig:7.9}{The region under \(y = 2x^{2} - x^{3}\) is awkward to describe with slices perpendicular to the axis of rotation: a horizontal strip’s two ends lie on the curve, so reading its radii means solving the cubic for \(x\). A vertical strip reads its height \(f(x)\) straight off the graph, the strip the shell method keeps.}
\end{figureblock}
\subsechead{The shell method}
Let \(f\) be continuous with \(f \ge 0\) on \([a, b]\), where \(0 \le a < b\), and rotate the region under \(y = f(x)\) about the \(y\)-axis. Cut \([a, b]\) into \(n\) strips of width \(\Delta x\) as usual. The strip over \([x_{i-1}, x_i]\), swept around the axis, produces a shell whose inner wall has radius \(x_{i-1}\) and whose outer wall has radius \(x_i\). As in §7.2, the modelling assumption comes first, and we state it explicitly. We take the solid’s volume to be the limit of the volumes of these shell stacks as the strips shrink. This is the same assumption, made for the same reason, as the slab model of Definition \ref{def:7.2}, only with the slices now nested around the axis instead of stacked along it.

What makes the shell decomposition so effective is a small algebraic accident. Approximate the shell over \([x_{i-1}, x_i]\) by an \emph{annular cylinder} of constant height \(f(\bar x_i)\), where \(\bar x_i = \tfrac{1}{2}(x_{i-1} + x_i)\) is the strip’s midpoint. (The approximation is in the height: the solid’s actual wall follows the varying graph, while the cylinder’s wall is flat.) For the annular cylinder, we can compute the volume exactly by subtracting the inner cylinder from the outer one. The midpoint then simplifies that exact volume:

\[ \pi x_i^{2}\,f(\bar x_i) - \pi x_{i-1}^{2}\,f(\bar x_i) = \pi\,(x_i + x_{i-1})(x_i - x_{i-1})\,f(\bar x_i) = 2\pi\,\bar x_i\,f(\bar x_i)\,\Delta x, \]

since \(x_i + x_{i-1} = 2\bar x_i\) and \(x_i - x_{i-1} = \Delta x\). Radius times height times thickness, with a factor \(2\pi\): the shell’s worth is \emph{circumference times height times thickness}, as if the shell had been slit open and unrolled into a flat slab of length \(2\pi \bar x_i\), height \(f(\bar x_i)\), and thickness \(\Delta x\).

\begin{figureblock}
  \includegraphics[width=75.93mm]{ch07/figs/strip-to-shell}
\figcaption{fig:7.10}{The vertical strip at distance \(x\) from the axis, swept into a thin-walled cylindrical shell: radius \(x\) read \emph{radially} from the axis, height \(f(x)\) read from the graph, thickness \(\Delta x\) between the walls. A shallower shell sits dotted inside. The solid is built of nested tubes, not stacked slices.}
\end{figureblock}
\begin{figureblock}
  \includegraphics[width=79.9mm]{ch07/figs/shell-unrolled}
\figcaption{fig:7.11}{The shell slit open and unrolled, \emph{as if} its curved wall lay flat: a slab of length \(2\pi\,\bar x_i\) (the circumference), height \(f(\bar x_i)\), and thickness \(\Delta x\). The picture explains where the \(2\pi\) comes from; the exact statement is the midpoint identity displayed above, not the flattening itself.}
\end{figureblock}
Summing the stack gives \(\sum_{i=1}^{n} 2\pi\,\bar x_i\,f(\bar x_i)\,\Delta x\), and this is precisely a Riemann sum for the function \(2\pi x\,f(x)\), sampled at the midpoints. Definition 6.2 permits \emph{any} choice of sample point, midpoints included; and since \(2\pi x f(x)\) is continuous on \([a, b]\), Theorem 6.1 guarantees the sums converge to its integral, no matter how the strips shrink. The shell stacks therefore have a limit, and the model hands that limit to the solid:

\begin{envtheorem}{Theorem}{thm:7.1}{The shell method}
Let \(f\) be continuous with \(f(x) \ge 0\) on \([a, b]\), where \(0 \le a < b\). The solid obtained by rotating the region under \(y = f(x)\), \(a \le x \le b\), about the \(y\)-axis has volume

\[ V = \int_{a}^{b} 2\pi x\,f(x)\,dx. \]
\end{envtheorem}
\begin{envcaution}{Caution}{}{}
A shell’s radius is the distance from the strip to the \emph{axis of rotation}. For rotation about the \(y\)-axis, that distance is the horizontal coordinate \(x\), even though the strip itself stands vertically. The height, meanwhile, is read from the graph as usual. Keep the two roles separate: radius says how far around the strip travels, height says how tall it is.
\end{envcaution}
\begin{workedexample}
\begin{envexample}{Example}{ex:7.10}{}
Find the volume of the solid obtained by rotating the region under \(y = 2x^{2} - x^{3}\), from \(x = 0\) to \(x = 2\), about the \(y\)-axis.
\end{envexample}
\begin{envsolution}{Solution}{}{}
This is the section’s opening problem. First check the theorem’s hypotheses: \(2x^{2} - x^{3} = x^{2}(2 - x)\) is continuous and non-negative on \([0, 2]\) (it vanishes at both ends and is positive between), and the interval sits in \(x \ge 0\). The shell at position \(x\) has radius \(x\) and height \(2x^{2} - x^{3}\), so Theorem \ref{thm:7.1} gives

\[ \begin{aligned}
          V &= \int_{0}^{2} 2\pi x\,\bigl(2x^{2} - x^{3}\bigr)\,dx = 2\pi \int_{0}^{2} \bigl(2x^{3} - x^{4}\bigr)\,dx \\[2pt]
            &= 2\pi \left[\frac{x^{4}}{2} - \frac{x^{5}}{5}\right]_{0}^{2} = 2\pi\left(8 - \frac{32}{5}\right) = \frac{16\pi}{5}.
        \end{aligned} \]

The cubic never had to be inverted: the shell integral reads \(f\) exactly as the region presents it. \qedmark
\end{envsolution}
\end{workedexample}
When the rotated region lies between \emph{two} graphs, the shell’s height is the gap between them. This gap is the same quantity that served as the integrand in §7.1. The solid is the difference of two under-a-graph solids, so by the linearity of the integral (Theorem 6.2) their shell integrals subtract, and the height \(f(x) - g(x)\) rides in place of \(f(x)\).

\begin{workedexample}
\begin{envexample}{Example}{ex:7.11}{}
Find the volume of the solid obtained by rotating the region enclosed by \(y = x\) and \(y = x^{2}\) about the \(y\)-axis.
\end{envexample}
\begin{envsolution}{Solution}{}{}
The region is the one from Example \ref{ex:7.7}: between \(x = 0\) and \(x = 1\), with \(y = x\) above and \(y = x^{2}\) below. The shell at position \(x\) has radius \(x\) and height \(x - x^{2}\), so

\[ V = 2\pi \int_{0}^{1} x\,\bigl(x - x^{2}\bigr)\,dx = 2\pi \int_{0}^{1} \bigl(x^{2} - x^{3}\bigr)\,dx = 2\pi\left(\frac{1}{3} - \frac{1}{4}\right) = \frac{\pi}{6}. \]

Note what changed from Example \ref{ex:7.7} and what did not. The region is the same, but the axis is different. There we rotated about the \(x\)-axis and sliced into washers; here the \(y\)-axis makes shells natural. Different axis, different solid, different volume. \qedmark
\end{envsolution}
\end{workedexample}
\subsechead{Washer or shell?}
Both methods compute volumes of revolution; the craft is in choosing. The washer slices \emph{perpendicular} to the axis of rotation and asks for the region’s boundaries as seen from that axis; the shell slices \emph{parallel} to the axis and asks for them as the region naturally presents them. One question settles it: which strip has its length written directly in the given functions?

\begin{envstrategy}{Strategy}{strat:7.2}{Washer or shell?}
\begin{steps}
  \item \emph{Draw both strips.} Sketch the region with one strip perpendicular to the axis of rotation (a washer slice) and one parallel to it (a shell slice).
  \item \emph{Pick the strip that reads the boundaries directly.} If a strip’s ends would require solving \(y = f(x)\) for \(x\) (or vice versa), the other strip is usually the better choice.
  \item \emph{Assemble the integrand.} Washer: \(\pi\bigl(R^{2} - r^{2}\bigr)\), radii measured from the axis to the far and near boundaries. Shell: \(2\pi\,(\text{radius})(\text{height})\), radius from the axis to the strip, height the boundary gap.
  \item \emph{Shifted axis.} If the axis of rotation is a line \(x = c\) rather than a coordinate axis, nothing changes except that every radius becomes a distance to that line, \(\lvert x - c\rvert\), provided the region is kept entirely on one side of the axis.
\end{steps}
\end{envstrategy}
The two methods had better agree when both apply. The next example runs one solid through both computations.

\begin{workedexample}
\begin{envexample}{Example}{ex:7.12}{}
The region under \(y = x^{2}\), from \(x = 0\) to \(x = 1\), is rotated about the \(y\)-axis. Find the volume of the resulting solid twice: once by shells, once by washers.
\end{envexample}
\begin{envsolution}{Solution}{}{}
\emph{Shells.} The shell at position \(x\) has radius \(x\) and height \(x^{2}\):

\[ V = 2\pi \int_{0}^{1} x\cdot x^{2}\,dx = 2\pi\cdot\frac{1}{4} = \frac{\pi}{2}. \]

\emph{Washers.} Slice horizontally at height \(y\), with \(0 \le y \le 1\). The region at that height runs from the parabola out to the line \(x = 1\); since \(x^{2}\) is increasing on \([0, 1]\), solving \(y = x^{2}\) for the boundary gives \(x = \sqrt{y}\), the description of the boundary as \(x\) in terms of \(y\) that §7.1 used for horizontal slicing. Rotated, the slice is a washer with outer radius \(R = 1\) and inner radius \(r = \sqrt{y}\):

\[ V = \int_{0}^{1} \pi\bigl(1^{2} - (\sqrt{y})^{2}\bigr)\,dy = \pi \int_{0}^{1} (1 - y)\,dy = \pi\left(1 - \frac{1}{2}\right) = \frac{\pi}{2}. \]

The two computations slice the same solid in perpendicular directions and meet at the same number, \(\tfrac{\pi}{2}\), as they must if both models measure the one volume. That this agreement holds in general, for every solid both methods reach, is a fact we are not yet equipped to prove economically; the tools arrive with integration by parts in Chapter 8 and change of variables in Chapter 15. For now, each such pair of computations is its own check. \qedmark
\end{envsolution}
\end{workedexample}
\begin{figureblock}
  \includegraphics[width=60.06mm]{ch07/figs/washer-vs-shell-1}\hspace{6mm}\includegraphics[width=60.06mm]{ch07/figs/washer-vs-shell-2}
\figcaption{fig:7.12}{One solid, two slicings. The solid of Example \ref{ex:7.12} is a cylinder with a paraboloid bowl carved from its top. A horizontal slice at height \(y\) is a washer from \(\sqrt{y}\) out to \(1\); a vertical slice at radius \(x\) is a shell of height \(x^{2}\). Both decompositions must — and do — measure the same volume.}
\end{figureblock}
The section’s addition to the toolkit is not a new integral but a new freedom: the direction of slicing is ours to choose, and Strategy \ref{strat:7.2} turns that choice into a two-line ritual. Underneath, the honesty is the same as ever. The modelling commitment is stated once (solids are measured by their shell stacks), each approximating piece is computed exactly, and Chapter 6’s convergence machinery closes the limit. So far the schema has measured geometry: areas, and now volumes three ways. The next section applies the schema to a quantity from physics that accumulates along a motion rather than across a region. There, for the first time, the slicing must decide \emph{what} to slice.


\sechead{7.4}{Work}
So far the slicing schema has measured geometry, first areas and then volumes. This section takes it somewhere new: to \emph{work}, the physicist’s measure of the effort a force expends in moving something. When a constant force pushes an object through a distance, the work done is simply force times distance. But most forces of interest are not constant. Stretching a spring, you feel the resistance grow: the final centimetre costs more effort than the first, so no single product \(\text{force} \times \text{distance}\) can price the whole stretch. The remedy is the one this chapter keeps applying. Over a \emph{thin slice} of the motion the force barely changes, the product is honest again, and summing the slices builds an integral.

We adopt the physics as our modelling input, the way §7.2 adopted the cylinder. When a constant force of magnitude \(F\) moves an object through a distance \(d\) along the force’s own direction, the \emph{work done} is \(W = Fd\). In SI units force is measured in newtons (N) and distance in metres, so work carries the unit newton-metre, called the \emph{joule} (J). (Imperial texts measure the same quantity in foot-pounds; we will stay with SI.) Now let the force vary along the line of motion. Say the object moves along the \(x\)-axis from \(a\) to \(b\), pushed by a force whose component along the motion at position \(x\) is \(F(x)\). We slice \([a, b]\) into strips of width \(\Delta x\), price each strip at \(F(x_i^{*})\,\Delta x\), and pass to the limit. The sums are Riemann sums of \(F\), so for continuous \(F\) Theorem 6.1 supplies the limit, and we take it as the definition.

\begin{envdefinition}{Definition}{def:7.3}{Work}
Suppose an object moves along the \(x\)-axis from \(x = a\) to \(x = b\), acted on by a force whose component along the direction of motion is \(F(x)\) at position \(x\), with \(F\) continuous on \([a, b]\). The \emph{work} done by the force is

\[ W = \int_{a}^{b} F(x)\,dx. \]
\end{envdefinition}
Read as in §6.4: work is the net change accumulated from its rate, and force is the rate at which work accrues per metre of progress. A force opposing the motion has negative component and contributes negative work. The sign convention is built into the word \emph{component}.

\subsechead{Springs}
The classic varying force is a spring’s. \emph{Hooke’s law} is an empirical model, accurate for stretches that are not too large. It says that the force required to hold a spring stretched \(x\) units beyond its natural length is proportional to \(x\):

\[ F(x) = kx, \]

where \(k\), the \emph{spring constant}, measures the spring’s stiffness in newtons per metre. Two readings of the same law are worth separating: \(kx\) is the force \emph{we} apply to hold the stretch, while the spring itself pulls back with the opposite force \(-kx\). Computing \(\int F\,dx\) with \(F = kx\) therefore prices the work done \emph{against} the spring, the effort of stretching it.

\begin{figureblock}
  \includegraphics[width=70.64mm]{ch07/figs/spring-force-area}
\figcaption{fig:7.13}{Hooke’s law drawn: the applied force \(F = kx\) against the stretch \(x\) beyond natural length. The work of stretching from \(a\) to \(b\) is the shaded area under the force line, accumulated strip by strip, and the strips grow taller toward \(b\), which is why the last centimetre costs the most.}
\end{figureblock}
\begin{workedexample}
\begin{envexample}{Example}{ex:7.13}{}
A force of \(30\) N is required to hold a spring stretched \(0.1\) m beyond its natural length. How much work is done in stretching it from \(0.1\) m to \(0.2\) m beyond its natural length?
\end{envexample}
\begin{envsolution}{Solution}{}{}
Here the slice is a slice of \emph{displacement}: the moving end of the spring advances from \(x = 0.1\) to \(x = 0.2\), and the force varies as it goes. First pin down the model. Hooke’s law with the given data reads \(30 = k\,(0.1)\), so \(k = 300\) N/m and \(F(x) = 300x\). Then

\[ W = \int_{0.1}^{0.2} 300x\,dx = 150\bigl[x^{2}\bigr]_{0.1}^{0.2} = 150\,(0.04 - 0.01) = 4.5 \text{ J}. \]

For comparison, the first \(0.1\) m of stretch, from the natural length, costs only \(150\,(0.01) = 1.5\) J, the same distance at a third of the price. The force grows with the stretch, and the work grows faster than the distance. \qedmark
\end{envsolution}
\end{workedexample}
\subsechead{Lifting and pumping}
In the spring problem one object moved and the force on it varied. The next two problems are built differently: a cable, or a tankful of water, is an extended body whose parts each move a \emph{different} distance. As with the spring, what we compute is work done \emph{against} a resisting force, here gravity. To raise a slice steadily, the winch or pump applies an upward force whose magnitude equals the slice’s weight, and it is \emph{that} applied force, pointing along the upward motion, whose work the integrals below price: (the slice’s weight) \(\times\) (the distance \emph{it} rises). (Gravity itself pulls against the motion and therefore does negative work; overcoming it is what makes lifting require effort.)

The force required on any one slice stays constant as it rises; what varies from slice to slice is the journey. So the slicing must follow the structure: slice the \emph{material itself}, and price each slice’s lift. Deciding what to slice determines the rest of the solution, and it is the first sentence of each solution below.

\begin{workedexample}
\begin{envexample}{Example}{ex:7.14}{}
A \(20\) m cable weighing \(2\) N per metre hangs straight down from a winch. How much work does the winch do in winding the whole cable up to the top?
\end{envexample}
\begin{envsolution}{Solution}{}{}
Slice the \emph{cable}. Measure \(x\) downward from the winch, and consider the slice of cable between depths \(x\) and \(x + \Delta x\). Its weight is (weight density) \(\times\) (length) \(= 2\,\Delta x\) newtons. The given \(2\) N/m is already a weight per metre, so no conversion from mass is needed. To reach the winch this slice must rise the distance \(x\) from its starting depth. Hauling it up steadily, the winch applies an upward force of that same magnitude, \(2\,\Delta x\) newtons, through the distance \(x\), contributing \(2x\,\Delta x\) joules of work against gravity; and the whole cable costs

\[ W = \int_{0}^{20} 2x\,dx = \bigl[x^{2}\bigr]_{0}^{20} = 400 \text{ J}. \]

The top of the cable rises almost no distance and contributes almost nothing; the bottom metre rises nearly the full \(20\) m and is the expensive part. The integral prices each slice by its own journey. \qedmark
\end{envsolution}
\end{workedexample}
\begin{figureblock}
  \includegraphics[width=49.48mm]{ch07/figs/cable-slice}
\figcaption{fig:7.14}{One slice of the hanging cable, at depth \(x\) below the winch. The winch carries its weight through the distance \(x\) that this slice alone must rise: slices near the bottom travel farthest and cost the most.}
\end{figureblock}
\begin{workedexample}
\begin{envexample}{Example}{ex:7.15}{}
A tank has the shape of an inverted circular cone (vertex at the bottom) of height \(4\) m and top radius \(2\) m. It is filled with water to a depth of \(3\) m. Find the work required to pump all the water out over the top rim. (Water weighs \(\rho g \approx 9800\) N per cubic metre: density \(\rho = 1000\) kg/m³ times \(g = 9.8\) m/s².)
\end{envexample}
\begin{envsolution}{Solution}{}{}
Slice the \emph{water}, horizontally. Each horizontal sheet of water rises as one piece, while sheets at different depths rise different amounts. Fix the coordinate first: let \(y\) run \emph{upward from the vertex}, so the tank occupies \(0 \le y \le 4\), the rim sits at \(y = 4\), and the water fills \(0 \le y \le 3\).

The slab of water at height \(y\), of thickness \(\Delta y\), is a disk. Its radius \(r(y)\) comes from similar triangles in the cone’s side view, exactly as in Example \ref{ex:7.9}: \(r(y)/2 = y/4\), so \(r(y) = y/2\). The slab’s volume is \(\pi\,(y/2)^{2}\,\Delta y = \tfrac{\pi}{4}y^{2}\,\Delta y\) cubic metres, so its weight is \(9800\cdot\tfrac{\pi}{4}y^{2}\,\Delta y\) newtons; raising it steadily to the rim takes an upward force of that magnitude acting from height \(y\) to height \(4\), through the lift distance \(4 - y\) metres. Summing the slabs’ work against gravity and refining,

\[ \begin{aligned}
          W &= \int_{0}^{3} 9800\cdot\frac{\pi}{4}\,y^{2}\,(4 - y)\,dy = 2450\pi \int_{0}^{3} \bigl(4y^{2} - y^{3}\bigr)\,dy \\[2pt]
            &= 2450\pi \left[\frac{4y^{3}}{3} - \frac{y^{4}}{4}\right]_{0}^{3} \\[2pt]
            &= 2450\pi\left(36 - \frac{81}{4}\right) = 2450\pi\cdot\frac{63}{4} \approx 1.2 \times 10^{5} \text{ J}.
        \end{aligned} \]

Every factor in the integrand earned its place from the slicing: \(\tfrac{\pi}{4}y^{2}\) is the slab’s cross-section (§7.2’s machinery on loan), \(9800\) converts volume to weight, and \(4 - y\) is that slab’s own lift. Note the two different numbers \(3\) and \(4\): the water surface bounds the \emph{integral}, while the rim height enters the \emph{distance}. \qedmark
\end{envsolution}
\end{workedexample}
\begin{figureblock}
  \includegraphics[width=70.64mm]{ch07/figs/tank-pump}
\figcaption{fig:7.15}{The conical tank in side view: vertex at height \(0\), rim at \(4\), water to depth \(3\). The slab at height \(y\) has radius \(y/2\) by similar triangles and rises \(4 - y\) to the rim. The water surface bounds the integral at \(3\), while the rim height \(4\) enters each slab’s distance.}
\end{figureblock}
\begin{envcaution}{Caution}{}{}
The three examples hide two different slicings. The spring slices the \emph{displacement}: one object moves, the applied force on it varies, and \(dx\) is a bit of distance. The cable and the tank slice the \emph{material}: the force required on each slice is constant (matching its weight), and what varies is the distance that slice travels. Memorizing a single template — “weight times height” or “force times \(dx\)” — will misprice one of the two. Name what is being sliced first; the integrand then assembles itself.
\end{envcaution}
\begin{envstrategy}{Strategy}{strat:7.3}{Setting up a work integral}
\begin{steps}
  \item \emph{Name the slice.} Decide what is being cut into thin pieces: the displacement (a varying force on one moving object) or the material itself (each piece moving its own distance).
  \item \emph{Price one slice.} Write the force on the slice and the distance that slice moves, with units attached to each factor.
  \item \emph{Assemble the contribution.} The slice’s work is (force) \(\times\) (distance); check that its units come to newton-metres, that is, joules.
  \item \emph{Integrate over the slices.} The limits of integration describe where the slices \emph{are} (the stretch traversed, the cable’s depths, the water’s heights), not necessarily where they end up.
\end{steps}
\end{envstrategy}
The schema has now left geometry and priced a motion. Nothing in the machinery changed. The sums were Riemann sums as ever, and Theorem 6.1 closed each limit. But for the first time the slicing had to make a decision about \emph{what} to slice, and Strategy \ref{strat:7.3} begins there. The next section asks the integral a different kind of question. Instead of accumulating a total — an area, a volume, a number of joules — it will extract a single representative number from a continuously varying function: its average.


\sechead{7.5}{Average Value of a Function}
The average of finitely many numbers is a familiar recipe: add them, divide by how many there are. Now ask for the average temperature over a day. Temperature varies continuously. There are not finitely many values to add, and no “how many” to divide by. The natural move is to sample: read the thermometer at \(n\) evenly spaced moments and average the readings. Different sampling rates give slightly different answers, but as the sampling grows finer the answers settle toward a single number, and that settling is by now a familiar sight. Write it out. If we sample \(f\) at the points \(x_1, \dots, x_n\) of a regular partition of \([a, b]\), with spacing \(\Delta x = \tfrac{b-a}{n}\), the sampled mean is

\[ \frac{f(x_1) + \dots + f(x_n)}{n} = \frac{1}{b - a}\sum_{i=1}^{n} f(x_i)\,\Delta x, \]

since \(\tfrac{1}{n} = \tfrac{\Delta x}{b-a}\). The sampled mean \emph{is} a Riemann sum, scaled by the interval’s width. So for continuous \(f\), Theorem 6.1 says the means converge, and their limit is an integral. The integral plays the role of “the sum of all the values”, and the width \(b - a\) plays the role of “how many”.

\begin{envdefinition}{Definition}{def:7.4}{Average value of a function}
Let \(f\) be continuous on \([a, b]\) with \(a < b\). The \emph{average value} of \(f\) on \([a, b]\) is

\[ f_{\text{ave}} = \frac{1}{b - a} \int_{a}^{b} f(x)\,dx. \]
\end{envdefinition}
The geometry behind the formula is worth keeping in view. Rewrite it as \(f_{\text{ave}}\,(b - a) = \int_{a}^{b} f(x)\,dx\): the rectangle over \([a, b]\) with height \(f_{\text{ave}}\) has the same (signed) area as the region under \(f\). The average is the \emph{levelling height}. Pour the region under the graph into a rectangular mould on the same base, and \(f_{\text{ave}}\) is where the surface comes to rest.

\begin{figureblock}
  \includegraphics[width=71.97mm]{ch07/figs/levelling-rectangle}
\figcaption{fig:7.16}{Levelling the region under \(y = 1 + x^{2}\) on \([-1, 2]\): the rectangle of height \(f_{\text{ave}}\) holds the same area as the region under the curve. The deficit above the curve on \([-1, 1]\) exactly balances the surplus above the lid on \([1, 2]\), so the levelling rectangle and the region hold the same area.}
\end{figureblock}
One everyday average is already an instance of this. Let \(f = v\) be the velocity of an object moving along a line; then its average value over \([a, b]\) is

\[ v_{\text{ave}} = \frac{1}{b - a}\int_{a}^{b} v(t)\,dt = \frac{s(b) - s(a)}{b - a}, \]

since the integral of velocity is the net change in position \(s\), by the Net Change Theorem of §6.4. That last ratio is total displacement over elapsed time: the \emph{average velocity} of elementary kinematics. The levelling-height rectangle is that familiar average, read now for any continuously varying quantity.

\begin{envcaution}{Caution}{}{}
The average value belongs to the function and the interval \emph{together}, not to the function alone. “The average of \(x^{2}\)” means nothing until the interval is named, and the same function averages differently over different stretches, as the next two examples show side by side.
\end{envcaution}
\begin{workedexample}
\begin{envexample}{Example}{ex:7.16}{}
Find the average value of \(f(x) = 1 + x^{2}\) on the interval \([-1, 2]\).
\end{envexample}
\begin{envsolution}{Solution}{}{}
Here \(a = -1\) and \(b = 2\), so \(b - a = 3\) and

\[ \begin{aligned}
          f_{\text{ave}} &= \frac{1}{3} \int_{-1}^{2} \bigl(1 + x^{2}\bigr)\,dx = \frac{1}{3}\left[x + \frac{x^{3}}{3}\right]_{-1}^{2} \\[2pt]
            &= \frac{1}{3}\left(\left(2 + \frac{8}{3}\right) - \left(-1 - \frac{1}{3}\right)\right) = \frac{1}{3}\cdot 6 = 2.
        \end{aligned} \]

The values of \(f\) on \([-1, 2]\) run from \(1\) up to \(5\), and their levelling height is \(2\). That is closer to the low end, because the function spends most of the interval below its peak. \qedmark
\end{envsolution}
\end{workedexample}
\begin{workedexample}
\begin{envexample}{Example}{ex:7.17}{}
Find the average value of the same function, \(f(x) = 1 + x^{2}\), on the interval \([0, 1]\).
\end{envexample}
\begin{envsolution}{Solution}{}{}
Now \(b - a = 1\), and

\[ f_{\text{ave}} = \int_{0}^{1} \bigl(1 + x^{2}\bigr)\,dx = \left[x + \frac{x^{3}}{3}\right]_{0}^{1} = \frac{4}{3}. \]

Same function, different interval, different average: on \([0, 1]\) the function never rises above \(2\), and its levelling height drops to \(\tfrac{4}{3}\). The average is a property of \(f\) \emph{together with} the interval, precisely the caution above. \qedmark
\end{envsolution}
\end{workedexample}
\subsechead{The Mean Value Theorem for Integrals}
In Example \ref{ex:7.16} the average was \(2\), and the function actually \emph{takes} the value \(2\) at \(x = \pm 1\), where the curve crosses the top of its levelling rectangle. Is that a coincidence? For a finite list it can easily fail: the average of the two numbers \(1\) and \(3\) is \(2\), which is neither of them. For a continuous function this cannot happen. Its values form an unbroken range, and the average lies between the smallest and the largest of them. So the average must itself be a value the function passes through. That is the content of the next theorem, and the machinery to prove it is exactly the pair of existence theorems Chapter 4 established for continuous functions on closed intervals.

\begin{envtheorem}{Theorem}{thm:7.2}{Mean Value Theorem for Integrals}
Let \(f\) be continuous on \([a, b]\) with \(a < b\). Then there is a point \(c\) in \([a, b]\) at which \(f\) takes its average value:

\[ f(c) = f_{\text{ave}} = \frac{1}{b - a}\int_{a}^{b} f(x)\,dx. \]
\end{envtheorem}
\begin{envproof}{Proof}{}{}
By the Extreme Value Theorem (Theorem 4.9(a)), \(f\) attains a minimum value \(m = f(x_m)\) and a maximum value \(M = f(x_M)\) at some points \(x_m, x_M\) of \([a, b]\). Since \(m \le f(x) \le M\) on \([a, b]\), the comparison property of the integral (Theorem 6.2, part 4) gives

\[ m\,(b - a) \le \int_{a}^{b} f(x)\,dx \le M\,(b - a), \]

and dividing by the positive number \(b - a\) traps the average: \(m \le f_{\text{ave}} \le M\).

If \(m = M\), the trap is already closed: \(f\) is constant on \([a, b]\), every value equals \(f_{\text{ave}}\), and any \(c\) in \([a, b]\) serves.

If \(m < M\), let \(\alpha\) be the smaller of \(x_m, x_M\) and \(\beta\) the larger. Then \([\alpha, \beta]\) is a genuine subinterval of \([a, b]\) whose endpoints carry the values \(m\) and \(M\) in one order or the other. Passing to this subinterval is what makes the next step work: the Intermediate Value Theorem needs an interval across which \(f\) runs from \(m\) to \(M\), and the endpoints \(a, b\) need not take those extreme values, whereas \(\alpha, \beta\) do. Since \(f\) is continuous on \([\alpha, \beta]\) and \(f_{\text{ave}}\) lies between its endpoint values, the Intermediate Value Theorem (Theorem 4.9(b)) supplies a point \(c\) in \([\alpha, \beta]\) with \(f(c) = f_{\text{ave}}\). Since \([\alpha, \beta]\) lies inside \([a, b]\), this point also lies in \([a, b]\). \qedmark
\end{envproof}
We note, without carrying it out, a second route to the same theorem: applying the Mean Value Theorem for derivatives (Theorem 4.12) to the accumulation function \(\int_{a}^{x} f(t)\,dt\), whose derivative is \(f\) by the Fundamental Theorem (Theorem 6.3), shows that the two mean value theorems are two statements of one fact.

\begin{workedexample}
\begin{envexample}{Example}{ex:7.18}{}
For \(f(x) = 1 + x^{2}\) on \([-1, 2]\), find every point \(c\) whose existence Theorem \ref{thm:7.2} guarantees.
\end{envexample}
\begin{envsolution}{Solution}{}{}
Example \ref{ex:7.16} computed \(f_{\text{ave}} = 2\). We solve \(f(c) = 2\):

\[ 1 + c^{2} = 2 \quad\Longleftrightarrow\quad c^{2} = 1 \quad\Longleftrightarrow\quad c = \pm 1, \]

and both candidates lie in \([-1, 2]\). The theorem promises \emph{at least one} such point and says nothing about uniqueness. Here there are two, one at the interval’s left endpoint and one interior. Geometrically, they are the two crossings between the curve and the top of its levelling rectangle: the rectangle’s lid cuts the graph wherever the function passes through its own average. \qedmark
\end{envsolution}
\end{workedexample}
This section added one definition and one theorem, both small and both permanent. The average value \(\tfrac{1}{b-a}\int_a^b f\) is the honest limit of sampled means, and the Mean Value Theorem for Integrals says a continuous function must pass through its own average. It is the first result in this book that depends on our having proved the Extreme and Intermediate Value Theorems rather than assuming them. Both results return when integration moves to higher dimensions, where regions replace intervals and the same questions arise again. The remaining two sections return the chapter to geometry, with the subtlest slicing yet: measuring the \emph{length} of a curve, where the slices are not rectangles at all.


\sechead{7.6}{Arc Length}
What is the \emph{length} of a curve? For segments of straight lines the distance formula answers at once; for anything curved, it is worth admitting that we have no formula at all, not a poor one, none. Physically the answer seems clear: lay a string along the curve, pull it straight, measure it. The mathematical translation of the string is the chapter’s schema with a new kind of slice. Mark points along the curve and join them by \emph{chords}: the resulting polygonal path can be measured exactly, one distance formula per chord. The polygon replaces each curved arc by a straight chord, so it under-measures every arc it spans. But as the marked points crowd together, the chords hug the curve ever more closely. For the first time this chapter, the slice is not a rectangle: it is a small straight segment standing in for a small piece of bend.

Make it precise for the graph of a function. Let \(f\) be continuous on \([a, b]\) with \(a < b\), and divide \([a, b]\) by the usual regular partition \(x_0 = a, x_1, \dots, x_n = b\) with spacing \(\Delta x\). The partition marks the points \(P_i = \bigl(x_i, f(x_i)\bigr)\) on the curve, and the inscribed polygon \(P_0 P_1 \cdots P_n\) has length \(\sum_{i=1}^{n} \lvert P_{i-1}P_i\rvert\), a number we can compute for every \(n\). Refining, we define:

\begin{envdefinition}{Definition}{def:7.5}{Arc length}
Let \(f\) be continuous on \([a, b]\), \(a < b\), and let \(P_i = \bigl(x_i, f(x_i)\bigr)\) be the points of the graph over the regular partition of \([a, b]\) into \(n\) pieces. The \emph{arc length} of the graph of \(f\) from \(P_0\) to \(P_n\) is

\[ L = \lim_{n \to \infty} \sum_{i=1}^{n} \bigl\lvert P_{i-1}P_i \bigr\rvert, \]

provided this limit exists as a finite number.
\end{envdefinition}
Two housekeeping notes. Length is a size, not a signed quantity: no orientation convention accompanies it, and a degenerate “curve” over a single point has length \(0\). And the definition makes the same kind of genuine commitment as Definition \ref{def:7.2} and the forthcoming Definition \ref{def:7.6}. It declares the polygon limit to be what length \emph{means} for a graph, with the clause “provided the limit exists” carrying the honesty.

\begin{figureblock}
  \includegraphics[width=62.71mm]{ch07/figs/polygon-refinement-1}\hspace{6mm}\includegraphics[width=62.71mm]{ch07/figs/polygon-refinement-2}
\figcaption{fig:7.17}{Inscribed polygons hugging the curve. With the same horizontal spacing, chords come out long where the curve is steep and short where it is flat; doubling \(n\) keeps the old corners and visibly tightens the fit. The polygon under-measures every arc it cuts.}
\end{figureblock}
\subsechead{The arc length formula}
When does the limit exist, and what is it? Each chord’s length is governed by how steeply the graph climbs across its strip, so the natural hypothesis is control over the \emph{slope}. We say \(f\) is \emph{smooth on \([a, b]\)} if \(f\) is differentiable on \([a, b]\) and its derivative \(f'\) is continuous there, the derivatives at the two endpoints being the one-sided ones. In the standard notation, this is written \(f \in C^{1}([a, b])\). Smoothness is exactly what the next theorem needs: it will hand each chord to the Mean Value Theorem and then hand the resulting sum to Chapter 6.

\begin{envtheorem}{Theorem}{thm:7.3}{Arc Length Formula}
Let \(f\) be smooth on \([a, b]\), \(a < b\). Then the arc length of the graph of \(f\) over \([a, b]\) exists, and

\[ L = \int_{a}^{b} \sqrt{1 + \bigl[f'(x)\bigr]^{2}}\,dx. \]
\end{envtheorem}
\begin{envproof}{Proof}{}{}
Write \(\Delta y_i = f(x_i) - f(x_{i-1})\) for the rise across the \(i\)-th strip. By the distance formula, the \(i\)-th chord has length

\[ \bigl\lvert P_{i-1}P_i \bigr\rvert = \sqrt{(\Delta x)^{2} + (\Delta y_i)^{2}}. \]

The blanket hypothesis pays out locally: since \(f'\) exists on all of \([a, b]\), the function \(f\) is continuous on each closed strip \([x_{i-1}, x_i]\) and differentiable on its interior, which is precisely what the Mean Value Theorem (Theorem 4.12) requires there. It supplies a point \(x_i^{*}\) in \((x_{i-1}, x_i)\) with

\[ \Delta y_i = f'(x_i^{*})\,\Delta x. \]

Substituting and pulling the positive factor \(\Delta x\) out of the square root,

\[ \bigl\lvert P_{i-1}P_i \bigr\rvert = \sqrt{(\Delta x)^{2} + \bigl[f'(x_i^{*})\bigr]^{2}(\Delta x)^{2}} = \sqrt{1 + \bigl[f'(x_i^{*})\bigr]^{2}}\;\Delta x, \]

so the whole polygon has length

\[ \sum_{i=1}^{n} \bigl\lvert P_{i-1}P_i \bigr\rvert = \sum_{i=1}^{n} \sqrt{1 + \bigl[f'(x_i^{*})\bigr]^{2}}\;\Delta x. \]

This is a Riemann sum for the function \(g(x) = \sqrt{1 + [f'(x)]^{2}}\), with sample points we did not choose ourselves but received from the Mean Value Theorem. That is where the second half of the hypothesis is needed: \(f'\) is continuous on \([a, b]\), so \(g\) is continuous there, and Theorem 6.1 guarantees that the Riemann sums of \(g\) converge to \(\int_{a}^{b} g(x)\,dx\) for \emph{every} choice of sample points (Definition 6.2), including the points supplied by the Mean Value Theorem. The polygon lengths therefore converge, and their limit is the integral claimed. \qedmark
\end{envproof}
\begin{figureblock}
  \includegraphics[width=73.29mm]{ch07/figs/mvt-chord-strip}
\figcaption{fig:7.18}{One strip of the proof of Theorem \ref{thm:7.3}. The chord is the hypotenuse over the legs \(\Delta x\) and \(\Delta y_i\), and the Mean Value Theorem supplies the interior point \(x_i^{*}\) whose tangent runs parallel to the chord, the picture of Theorem 4.12, now put to work strip by strip.}
\end{figureblock}
The integrand deserves a reading of its own. Where the graph runs flat, \(f' = 0\) and \(\sqrt{1 + f'^{2}} = 1\): arc length accrues at the same rate as horizontal distance. Where the graph is steep, the factor grows, and the curve packs more length above each unit of run. The integrand is the local \emph{stretch factor} converting run into length, never less than \(1\), and equal to \(1\) only where the curve is momentarily level.

\begin{workedexample}
\begin{envexample}{Example}{ex:7.19}{}
Find the length of the curve \(y = x^{3/2}\) from \((0, 0)\) to \((1, 1)\).
\end{envexample}
\begin{envsolution}{Solution}{}{}
First, check the hypothesis. Here \(f(x) = x^{3/2}\), so \(f'(x) = \tfrac{3}{2}x^{1/2}\), which is continuous on all of \([0, 1]\). At the left endpoint the one-sided derivative is \(0\), and \(f'\) approaches it continuously. So \(f\) is smooth there and Theorem \ref{thm:7.3} applies:

\[ L = \int_{0}^{1} \sqrt{1 + \tfrac{9}{4}x}\;dx. \]

Substitute \(u = 1 + \tfrac{9}{4}x\), so \(du = \tfrac{9}{4}\,dx\) and \(dx = \tfrac{4}{9}\,du\), with limits \(u = 1\) to \(u = \tfrac{13}{4}\):

\[ \begin{aligned}
          L &= \frac{4}{9}\int_{1}^{13/4} \sqrt{u}\;du = \frac{4}{9}\cdot\frac{2}{3}\Bigl[u^{3/2}\Bigr]_{1}^{13/4} \\[2pt]
            &= \frac{8}{27}\left(\Bigl(\frac{13}{4}\Bigr)^{3/2} - 1\right) = \frac{13\sqrt{13} - 8}{27} \approx 1.44.
        \end{aligned} \]

A quick sanity check: the straight chord from \((0,0)\) to \((1,1)\) has length \(\sqrt{2} \approx 1.41\), and the curve, which bends away from the chord, is a little longer, as an inscribed polygon with one chord would already predict. \qedmark
\end{envsolution}
\end{workedexample}
The curve just measured is no arbitrary choice. The graph of \(y = x^{3/2}\) is the \emph{semicubical parabola} \(y^{2} = x^{3}\), and in 1657 the nineteen-year-old William Neile made it the first algebraic curve ever to have its length computed, a result his teacher Wallis soon published. The feat carried weight because Descartes, in his \emph{Géométrie}, had judged that the ratio of a curved line to a straight one “cannot be discovered by human minds”; Neile found one such length exactly, the first crack in that verdict.

\begin{workedexample}
\begin{envexample}{Example}{ex:7.20}{}
Find the length of the curve \(y = \dfrac{x^{3}}{6} + \dfrac{1}{2x}\) from \(x = 1\) to \(x = 2\).
\end{envexample}
\begin{envsolution}{Solution}{}{}
On \([1, 2]\) the function and its derivative \(f'(x) = \tfrac{x^{2}}{2} - \tfrac{1}{2x^{2}}\) are continuous (\(x = 0\) is safely outside the interval), so the formula applies. Squaring and adding \(1\) produces a perfect square:

\[ \begin{aligned}
          1 + \bigl[f'(x)\bigr]^{2} &= 1 + \frac{x^{4}}{4} - \frac{1}{2} + \frac{1}{4x^{4}} \\[2pt]
            &= \frac{x^{4}}{4} + \frac{1}{2} + \frac{1}{4x^{4}} = \left(\frac{x^{2}}{2} + \frac{1}{2x^{2}}\right)^{2}.
        \end{aligned} \]

Taking the square root demands a sign check: \(\sqrt{q^{2}} = \lvert q\rvert\), and here \(q = \tfrac{x^{2}}{2} + \tfrac{1}{2x^{2}}\) is a sum of two positive terms on \([1, 2]\), so the absolute value opens without a sign change. Hence

\[ \begin{aligned}
          L &= \int_{1}^{2} \left(\frac{x^{2}}{2} + \frac{1}{2x^{2}}\right) dx = \left[\frac{x^{3}}{6} - \frac{1}{2x}\right]_{1}^{2} \\[2pt]
            &= \left(\frac{8}{6} - \frac{1}{4}\right) - \left(\frac{1}{6} - \frac{1}{2}\right) = \frac{17}{12}.
        \end{aligned} \]

The curve was engineered so that \(1 + f'^{2}\) assembles into a square, which is why the answer is exact and tidy. For a curve chosen at random, \(\sqrt{1 + f'^{2}}\) rarely has an antiderivative among the functions we know; Chapter 8 will widen the toolkit, and where even that fails, its numerical methods take over. \qedmark
\end{envsolution}
\end{workedexample}
\begin{envstrategy}{Strategy}{strat:7.4}{Setting up an arc length integral}
\begin{steps}
  \item \emph{Check smoothness.} Confirm \(f'\) is continuous on \([a, b]\); at a corner, split the curve and add the lengths of the pieces.
  \item \emph{Form the integrand.} Write \(1 + \bigl[f'(x)\bigr]^{2}\) and simplify.
  \item \emph{Find structure.} Look for a structure that makes the root tractable, such as a substitution (as in Example \ref{ex:7.19}) or a perfect square (Example \ref{ex:7.20}); most other integrands wait for Chapter 8.
  \item \emph{If it is a perfect square, open the root carefully.} Writing \(1 + \bigl[f'(x)\bigr]^{2} = q^{2}\), the root is \(\sqrt{q^{2}} = \lvert q\rvert\); confirm the sign of \(q\) on \([a, b]\) before dropping the bars.
  \item \emph{Integrate.} Evaluate \(\int_{a}^{b}\sqrt{1 + [f'(x)]^{2}}\,dx\); where no elementary antiderivative exists, the length stands as a definite integral for the methods of Chapter 8.
\end{steps}
\end{envstrategy}
\subsechead{The arc length function and the differential \(ds\)}
Often what is wanted is not one length but length \emph{as a running quantity}, that is, how far along the curve we have travelled by the time we reach \(x\). For \(f\) smooth on \([a, b]\), define the \emph{arc length function}

\[ s(x) = \int_{a}^{x} \sqrt{1 + \bigl[f'(t)\bigr]^{2}}\;dt, \qquad a \le x \le b, \]

which measures the graph from its left end to the point above \(x\); thus \(s(a) = 0\) and \(s(b) = L\). Its integrand is continuous, so Part 1 of the Fundamental Theorem (Theorem 6.3) differentiates it on sight:

\[ \frac{ds}{dx} = \sqrt{1 + \left(\frac{dy}{dx}\right)^{2}} \;\ge\; 1. \]

\begin{figureblock}
  \includegraphics[width=42.42mm]{ch07/figs/ds-triangle}
\figcaption{fig:7.19}{The differential triangle: \(ds^{2} = dx^{2} + dy^{2}\). A memory device for the proved formula of Theorem \ref{thm:7.3}: the chord picture shrunk to a point. It is not a derivation of that formula.}
\end{figureblock}
The inequality repeats the stretch-factor reading. The curve always accumulates at least as fast as the run beneath it. In the language of differentials from §5.3, the same equation reads

\[ ds = \sqrt{1 + \left(\frac{dy}{dx}\right)^{2}}\;dx, \]

and squaring it gives the easily remembered \emph{\(ds^{2} = dx^{2} + dy^{2}\)}, an infinitesimal Pythagorean theorem that shrinks the proof’s chord picture to a point. Its status is worth being precise about: it is a mnemonic \emph{for} the proved formula, not a derivation of it. The honest argument is the one Theorem \ref{thm:7.3} already made. As a mnemonic it is excellent, and \(ds\) itself will be the working currency of the next section and of the later chapters that measure along curves.

\begin{workedexample}
\begin{envexample}{Example}{ex:7.21}{}
Find the arc length function for the curve \(y = \dfrac{x^{2}}{8} - \ln x\), measured from the point \((1, \tfrac{1}{8})\), for \(x \ge 1\).
\end{envexample}
\begin{envsolution}{Solution}{}{}
Here \(f'(x) = \tfrac{x}{4} - \tfrac{1}{x}\), continuous for \(x > 0\), so \(f\) is smooth on every interval \([1, x]\). The perfect-square pattern of Example \ref{ex:7.20} repeats:

\[ 1 + \bigl[f'(t)\bigr]^{2} = 1 + \frac{t^{2}}{16} - \frac{1}{2} + \frac{1}{t^{2}} = \left(\frac{t}{4} + \frac{1}{t}\right)^{2}, \]

and the bracket \(\tfrac{t}{4} + \tfrac{1}{t}\) is positive for \(t \ge 1\), so the root opens cleanly. Here \(t > 0\), so the absolute value in \(\int \tfrac{1}{t}\,dt = \ln\lvert t\rvert\) from §6.4 is simply \(t\). Hence

\[ s(x) = \int_{1}^{x} \left(\frac{t}{4} + \frac{1}{t}\right) dt = \left[\frac{t^{2}}{8} + \ln t\right]_{1}^{x} = \frac{x^{2} - 1}{8} + \ln x. \]

The checks confirm it: \(s(1) = 0\), as a length measured from the starting point must be; and differentiating returns \(s'(x) = \tfrac{x}{4} + \tfrac{1}{x} = \sqrt{1 + [f'(x)]^{2}}\), the Fundamental Theorem closing the loop. For instance, \(s(2) = \tfrac{3}{8} + \ln 2 \approx 1.07\): by \(x = 2\), just over one unit of curve has unspooled. \qedmark
\end{envsolution}
\end{workedexample}
\begin{envcaution}{Caution}{}{}
Smoothness is the hypothesis of the \emph{formula}, not of length itself. The graph of \(\lvert x\rvert\) on \([-1, 1]\) has a perfectly good length. Definition \ref{def:7.5} measures it, and splitting at the corner gives \(2\sqrt{2}\). But Theorem \ref{thm:7.3} cannot be applied across \(x = 0\), where no derivative exists for the Mean Value Theorem to use. At a corner, split the curve and add the pieces.
\end{envcaution}
The section ran the chapter’s schema at full honesty: a definition that says what length means (chords, refined), a theorem that computes it (the Mean Value Theorem converting each chord, Chapter 6 closing the limit), and a pair of working tools, the arc length function \(s\) and its differential \(ds\). Everything here concerns the graph of a function of \(x\); curves given by parametric equations must wait for Chapter 10, and curves through space for Chapter 13, where \(ds\) will reappear each time in suitably generalized dress. The next section puts \(ds\) straight to work: rotate a smooth curve about an axis, and ask not for the volume enclosed but for the \emph{area of the surface} swept out.


\sechead{7.7}{Area of a Surface of Revolution}
Rotating a region about an axis gave us solids and their volumes. Rotate just the \emph{curve}, and what it sweeps out is a surface rather than a solid, the skin with none of the filling. The natural question is its area. The slicing schema will serve once more, but the slice is the subtlest of the chapter. A curved surface cannot be flattened onto a plane (every mapmaker’s complaint about the globe), yet a thin \emph{band} of it very nearly can. Cut the surface into narrow bands, flatten each, add up. Note what a band involves that no earlier slice did: it is long one way — all the way around the axis — and narrow the other way, along the curve. Two length scales, circumference and slant, and the integrand will have to carry both.

\subsechead{Cones, frustums, and bands}
The flattenable model piece is a cone. Slit a cone of base radius \(r\) and slant height \(\ell\) down its side and unroll it: it lies flat as a circular sector of radius \(\ell\) whose arc has length \(2\pi r\), the rim of the base. By the sector-area formula, its area is \(\tfrac{1}{2}(\text{arc})(\text{radius}) = \tfrac{1}{2}(2\pi r)\,\ell = \pi r \ell\). That is the cone’s lateral area.

Our bands are not full cones but \emph{frustums}, which are cones with their tops cut off. Take a frustum of slant height \(\ell\) with end radii \(0 < r_1 < r_2\). Extending its side to the apex produces two cones with a common vertex, of slant heights \(\ell_1 < \ell_2\) with \(\ell_2 - \ell_1 = \ell\). The shared apex makes the two triangles similar, so \(\tfrac{\ell_1}{r_1} = \tfrac{\ell_2}{r_2}\). Hence \(\ell_2 = \tfrac{r_2}{r_1}\ell_1\). Substituting this into \(\ell_2 - \ell_1 = \ell\) gives \(\ell_1 = \tfrac{r_1\,\ell}{r_2 - r_1}\), and then \(\ell_2 = \tfrac{r_2\,\ell}{r_2 - r_1}\). Subtracting lateral areas,

\[ \pi r_2 \ell_2 - \pi r_1 \ell_1 = \pi\,\frac{r_2^{2} - r_1^{2}}{r_2 - r_1}\,\ell = \pi\,(r_1 + r_2)\,\ell. \]

The derivation assumed \(0 < r_1 < r_2\), because it divided by \(r_1\) to relate the two slant heights and by \(r_2 - r_1\) at the end. The two excluded cases can be checked directly. Equal radii make the frustum a cylindrical band, which unrolls into a flat rectangle of length \(2\pi r\) and width \(\ell\), with area \(2\pi r \ell = \pi\,(r + r)\,\ell\). And \(r_1 = 0\) makes it the cone itself, whose lateral area \(\pi r_2 \ell\) is what \(\pi\,(r_1 + r_2)\,\ell\) gives when \(r_1 = 0\). The same formula covers all three cases, and it is worth compressing: writing \(\bar r = \tfrac{1}{2}(r_1 + r_2)\) for the \emph{average radius},

\[ \text{lateral area of a frustum} = 2\pi\,\bar r\,\ell. \]

Average radius around, slant height along: the frustum’s area is already “circumference times width”, read at the band’s middle.

\begin{figureblock}
  \includegraphics[width=77.26mm]{ch07/figs/frustum-unrolled}
\figcaption{fig:7.20}{The frustum band slit open and unrolled: an annular sector, not a rectangle. Its inner and outer edges keep different lengths. The dashed midline arc has length \(2\pi \bar r\), and midline times width \(\ell\) is exactly the band’s area; only the cylinder band (\(r_1 = r_2\)) unrolls into a true rectangle. A picture of the formula the cones already proved, not a second derivation.}
\end{figureblock}
\begin{figureblock}
  \includegraphics[width=60.06mm]{ch07/figs/chord-to-band-1}\hspace{6mm}\includegraphics[width=60.06mm]{ch07/figs/chord-to-band-2}
\figcaption{fig:7.21}{Definition \ref{def:7.6}’s approximant made visible: one chord \(P_{i-1}P_i\) of the curve, and the frustum band it sweeps about the \(x\)-axis. The band’s slant edge \emph{is} the chord, not the arc. Its end radii are the endpoint heights \(f(x_{i-1})\) and \(f(x_i)\), and its slant length is \(\ell = \lvert P_{i-1}P_i\rvert\).}
\end{figureblock}
\subsechead{The surface area of a surface of revolution}
Now rotate the graph of a continuous \(f \ge 0\) on \([a, b]\), \(a < b\), about the \(x\)-axis. Mark the graph over the regular partition, as in §7.6, at the points \(P_i = \bigl(x_i, f(x_i)\bigr)\). Each chord \(P_{i-1}P_i\), swept around the axis, generates precisely a frustum band: end radii \(f(x_{i-1})\) and \(f(x_i)\), slant height \(\lvert P_{i-1}P_i\rvert\). The frustum formula prices each band exactly, and the sum of the bands approximates the surface. As with arc length, we let the definition say the rest.

\begin{envdefinition}{Definition}{def:7.6}{Surface area of revolution}
Let \(f\) be continuous with \(f \ge 0\) on \([a, b]\), \(a < b\), and let \(P_0, \dots, P_n\) be the points of its graph over the regular partition. The \emph{area} of the surface obtained by rotating the graph about the \(x\)-axis is

\[ S = \lim_{n \to \infty} \sum_{i=1}^{n} 2\pi\cdot\frac{f(x_{i-1}) + f(x_i)}{2}\cdot\bigl\lvert P_{i-1}P_i\bigr\rvert, \]

provided this limit exists as a finite number.
\end{envdefinition}
For smooth curves the limit not only exists but lands on an integral we can often evaluate, and this time the whole argument closes with tools already in hand.

\begin{envtheorem}{Theorem}{thm:7.4}{Surface Area Formula}
Let \(f\) be smooth on \([a, b]\) (§7.6) with \(f \ge 0\), \(a < b\). Then the surface area of the surface obtained by rotating the graph of \(f\) about the \(x\)-axis exists, and

\[ S = \int_{a}^{b} 2\pi\,f(x)\,\sqrt{1 + \bigl[f'(x)\bigr]^{2}}\;dx. \]
\end{envtheorem}
The proof that follows is the most technical argument in the chapter. On a first reading it is enough to accept Theorem \ref{thm:7.4} and turn to the examples; the derivation is here for completeness, and it rewards a second pass.

\begin{envproof}{Proof}{}{}
The chord conversion is already paid for: in the proof of Theorem \ref{thm:7.3}, the Mean Value Theorem gave, for each strip, a point \(\xi_i\) in \((x_{i-1}, x_i)\) with \(\lvert P_{i-1}P_i\rvert = \sqrt{1 + [f'(\xi_i)]^{2}}\,\Delta x\). The \(i\)-th band’s area is therefore

\[ 2\pi\cdot\frac{f(x_{i-1}) + f(x_i)}{2}\cdot\sqrt{1 + \bigl[f'(\xi_i)\bigr]^{2}}\;\Delta x, \]

which reads the curve’s height at the strip’s two \emph{ends} while reading its slope at the interior point \(\xi_i\). If the height were read at \(\xi_i\) as well, the band sum would be exactly a Riemann sum for \(2\pi f\sqrt{1 + f'^{2}}\); so we estimate the cost of moving the height reading, and show it vanishes in the limit.

Smoothness makes both \(\lvert f'\rvert\) and \(\sqrt{1 + f'^{2}}\) continuous on \([a, b]\), so by the Extreme Value Theorem (Theorem 4.9(a)) they attain maxima: say \(K = \max \lvert f'\rvert\) and \(B = \max \sqrt{1 + f'^{2}}\). For any point \(u\) in the \(i\)-th strip, the Mean Value Theorem gives \(f(u) - f(\xi_i) = f'(\zeta)\,(u - \xi_i)\) for some \(\zeta\) between them, and \(\lvert f'(\zeta)\rvert \le K\); so it bounds the height drift by the slope bound: \(\lvert f(u) - f(\xi_i)\rvert \le K\,\lvert u - \xi_i\rvert \le K\,\Delta x\). Applying this at \(u = x_{i-1}\) and \(u = x_i\), each of the two endpoint drifts is at most \(K\,\Delta x\), so by the triangle inequality their average is too:

\[ \left\lvert \frac{f(x_{i-1}) + f(x_i)}{2} - f(\xi_i) \right\rvert \le K\,\Delta x, \]

so replacing the \(i\)-th band’s height reading changes that term by at most \(2\pi\,(K \Delta x)\,B\,\Delta x\). Across all \(n\) bands the total change is at most

\[ n \cdot 2\pi K B\,(\Delta x)^{2} = 2\pi K B\,(b - a)\,\Delta x \;\longrightarrow\; 0 \qquad (n \to \infty), \]

since \(n\,\Delta x = b - a\). The corrected sums \(\sum 2\pi f(\xi_i)\sqrt{1 + [f'(\xi_i)]^{2}}\,\Delta x\) are Riemann sums of the continuous function \(2\pi f \sqrt{1 + f'^{2}}\), and by Theorem 6.1 with Definition 6.2’s any-sample-point clause they converge to its integral. The band sums differ from them by an amount tending to \(0\), so they converge to the same integral, which is the claimed formula. \qedmark
\end{envproof}
The integrand is worth a second look. Here \(2\pi f(x)\) is the circumference of the circle the point \((x, f(x))\) travels, and \(\sqrt{1 + f'^{2}}\,dx\) is the arc element \(ds\) of §7.6. These are the band’s two length scales, one around the axis and one along the curve, exactly as promised. In the differential shorthand the theorem compresses to

\[ S = \int 2\pi \rho\;ds, \]

where \(\rho\) is the distance from the moving point to the axis of rotation. About the \(x\)-axis, \(\rho = y = f(x)\). About the \emph{\(y\)}-axis, \(\rho = x\), with the curve kept entirely in \(x \ge 0\) (an interval \(0 \le a < b\)). Two things force that restriction: a signed \(x\) would count area negatively, and the two halves of a curve straddling the axis can overlap or re-cover parts of one surface. The general case therefore calls for splitting at the axis and tracking what each piece generates. One principle, two readings. The same MVT argument, run with the radius \(x\) in place of \(f(x)\), proves the \(y\)-axis version.

\begin{envcaution}{Caution}{}{}
Do not memorize \(2\pi y\,ds\) and \(2\pi x\,ds\) as two formulas. There is one: \(2\pi \rho\,ds\), with \(\rho\) the distance from the curve to the \emph{axis of rotation}. This is the same discipline that set every washer and shell radius in §§7.2–7.3. Identify the axis first; \(\rho\) follows.
\end{envcaution}
\begin{workedexample}
\begin{envexample}{Example}{ex:7.22}{}
A \emph{zone} of a sphere is the band cut from it by two parallel planes. Show that for the sphere of radius \(r\) generated by rotating \(y = \sqrt{r^{2} - x^{2}}\) about the \(x\)-axis, the zone between the planes \(x = a\) and \(x = b\) (where \(-r < a < b < r\)) has area \(2\pi r\,(b - a)\).
\end{envexample}
\begin{envsolution}{Solution}{}{}
On \([a, b]\), safely inside \((-r, r)\), the function \(f(x) = \sqrt{r^{2} - x^{2}}\) is positive and its derivative \(f'(x) = -x/\sqrt{r^{2} - x^{2}}\) is continuous. The curve is smooth there, and Theorem \ref{thm:7.4} applies. Compute the integrand:

\[ \begin{aligned}
          1 + \bigl[f'(x)\bigr]^{2} &= 1 + \frac{x^{2}}{r^{2} - x^{2}} = \frac{r^{2}}{r^{2} - x^{2}}, \\[2pt]
          \text{so}\qquad f(x)\sqrt{1 + \bigl[f'(x)\bigr]^{2}} &= \sqrt{r^{2} - x^{2}}\cdot\frac{r}{\sqrt{r^{2} - x^{2}}} = r.
        \end{aligned} \]

The heights and slopes conspire to cancel: where the sphere’s profile is low, it is steep, and the band’s area comes out the same everywhere. Hence

\[ S = \int_{a}^{b} 2\pi r\,dx = 2\pi r\,(b - a). \]

The zone’s area depends only on the \emph{width} \(b - a\) of the slab that cuts it, not on where along the sphere the cut is made. A band of given thickness near a pole has exactly the same area as one at the equator. \qedmark
\end{envsolution}
\end{workedexample}
This width-only law is ancient. Archimedes knew that a sphere and its circumscribed cylinder have equal areas in corresponding bands. If the same pair of parallel planes cuts both solids, the two bands have equal area. That correspondence is why the result is remembered as his \emph{hat-box theorem}. What the integral adds is the bookkeeping that makes the cancellation visible in two lines.

\begin{figureblock}
  \includegraphics[width=70.64mm]{ch07/figs/zone-bands}
\figcaption{fig:7.22}{Two slabs of the same width \(w\) cut the sphere’s profile, one at the equator and one out toward the pole, both strictly inside \((-r, r)\). The polar arc is longer and steeper but sits at a smaller radius; the equatorial arc is shorter and flatter at a larger one. Swept about the axis, the two arcs generate zones of \emph{equal} area: the compensation Example \ref{ex:7.22} computed.}
\end{figureblock}
Letting the two cutting planes recede to the poles, \(a \to -r\) and \(b \to r\), the zone areas approach \(2\pi r \cdot 2r = 4\pi r^{2}\), the classical area of the whole sphere. Honesty requires a flag here: at the poles themselves the profile has vertical tangents, \(f'\) blows up, and Theorem \ref{thm:7.4}’s smoothness hypothesis fails on the full interval \([-r, r]\), so the full-sphere value is reached as a \emph{limit of zones}, not by one application of the theorem. Taming endpoint blow-ups of exactly this kind is the business of improper integrals, in Chapter 8.

\begin{workedexample}
\begin{envexample}{Example}{ex:7.23}{}
The arc of the parabola \(y = x^{2}\) from \((0, 0)\) to \((1, 1)\) is rotated about the \(y\)-axis. Find the area of the resulting surface.
\end{envexample}
\begin{envsolution}{Solution}{}{}
The axis of rotation is the \(y\)-axis, so the radius is \(\rho = x\), and the arc lives in \(0 \le x \le 1\), on one side of the axis, as the \(y\)-axis form requires. With \(y = x^{2}\), the arc element is \(ds = \sqrt{1 + (2x)^{2}}\,dx = \sqrt{1 + 4x^{2}}\,dx\), so

\[ S = \int 2\pi\,\rho\;ds = 2\pi \int_{0}^{1} x\,\sqrt{1 + 4x^{2}}\;dx. \]

Substitute \(u = 1 + 4x^{2}\), \(du = 8x\,dx\), with limits \(u = 1\) to \(u = 5\):

\[ S = \frac{2\pi}{8}\int_{1}^{5} \sqrt{u}\;du = \frac{\pi}{4}\cdot\frac{2}{3}\Bigl[u^{3/2}\Bigr]_{1}^{5} = \frac{\pi}{6}\bigl(5\sqrt{5} - 1\bigr) \approx 5.33. \]

Note where each ingredient came from: the radius \(x\) from the axis (not from the curve’s formula), the element \(ds\) from §7.6, and the substitution from §6.5, three sections cooperating in four lines. \qedmark
\end{envsolution}
\end{workedexample}
\begin{figureblock}
  \includegraphics[width=66.67mm]{ch07/figs/dual-axis-radius}
\figcaption{fig:7.23}{One point of the arc \(y = x^{2}\), two radii: rotation about the \(x\)-axis reads the vertical distance \(f(x)\), rotation about the \(y\)-axis reads the horizontal distance \(x\). The radius is always the distance to the axis of rotation, one principle and not a memorized pair.}
\end{figureblock}
\subsechead{Chapter summary}
This chapter took the machine Chapter 6 built and taught it to measure. One schema ran through every section. Slice the quantity, approximate each slice by a product, sum, and refine. Each section supplied its own slice. Section 7.1 sliced regions into strips whose height is a \emph{gap}, defining the area between curves (Definition \ref{def:7.1}) as \(\int_a^b (f - g)\,dx\), with crossings handled by \(\int \lvert f - g\rvert\). Section 7.2 promoted the strip to a solid slab and volume to the integral of cross-sectional area (Definition \ref{def:7.2}), \(V = \int_a^b A(x)\,dx\), with disks and washers as its applications to solids of revolution. The method also recovers the classical volumes of the sphere and the pyramid. Section 7.3 sliced the other way, parallel to the axis, proving the shell method (Theorem \ref{thm:7.1}), \(V = \int_a^b 2\pi x f(x)\,dx\), and distilling the washer-or-shell decision into Strategy \ref{strat:7.2}.

Section 7.4 left geometry. The work done by a varying force is \(W = \int_a^b F(x)\,dx\) (Definition \ref{def:7.3}), where \(F(x)\) is the signed component of the force along the motion. The slicing also learned to choose what to slice, whether displacement, cable, or water. Section 7.5 turned totals into averages, \(f_{\text{ave}} = \frac{1}{b-a}\int_a^b f\,dx\) (Definition \ref{def:7.4}), and proved the Mean Value Theorem for Integrals (Theorem \ref{thm:7.2}): a continuous function passes through its own average. Section 7.6 measured the curve itself by refining chords (Definition \ref{def:7.5}) and proved the Arc Length Formula (Theorem \ref{thm:7.3}), \(L = \int_a^b \sqrt{1 + [f'(x)]^{2}}\,dx\), minting the arc element \(ds\); and §7.7 wrapped the chords into frustum bands to define surface area (Definition \ref{def:7.6}) and prove the Surface Area Formula (Theorem \ref{thm:7.4}), compressed as \(S = \int 2\pi \rho\,ds\).

Three theorems in the chapter were proved with the machinery of Chapter 4. The Mean Value Theorem for Integrals uses the Extreme and Intermediate Value Theorems, and the two geometric formulas use the Mean Value Theorem. This is the first extended payoff in the book from having proved those earlier results rather than assuming them.

What the chapter leaves owing is plain from its own examples: the integrals that measure lengths and surfaces strain the small table of antiderivatives we possess, and the curves here were chosen to fit the toolkit rather than the other way round. Chapter 8 pays that debt with a systematic account of integration techniques, and with improper integrals, which will also settle the question about the poles raised just after Example \ref{ex:7.22}. Farther ahead, \(ds\) returns for parametric curves in Chapter 10 and for space curves in Chapter 13, the average value and the slicing of volumes reappear over two-dimensional regions in Chapter 15, and surfaces return in their full generality in Chapter 16. The slicing schema itself, though, is now complete. It is the way integrals meet the world, and the rest of the book will use it routinely.
\end{document}
